\documentclass[10pt,twocolumn]{article}

\usepackage[utf8]{inputenc}
\usepackage[margin=0.75in]{geometry}
\usepackage{amsmath,amssymb,amsfonts,amsthm}
\usepackage{booktabs}
\usepackage{tabularx}
\usepackage{graphicx}
\usepackage{hyperref}
\usepackage{microtype}
\usepackage{cite}
\usepackage{physics}
\usepackage{xcolor}

\hypersetup{
    colorlinks=true,
    linkcolor=blue!70!black,
    citecolor=blue!70!black,
    urlcolor=blue!70!black
}

\newtheorem{theorem}{Theorem}
\newtheorem{lemma}[theorem]{Lemma}
\newtheorem{definition}{Definition}

\title{\textbf{Subspace-Constrained Quantum Portfolio Selection via Excitation-Preserving Mixers and Heavy-Hex SWAP Networks}}

\author{
    \textbf{Team Gizmo} \\
    Akhilesh Varma (\texttt{akhverm@gmail.com}) \quad Abhay Gupta (\texttt{guptasanjay6283@gmail.com}) \\
    \textit{IBM Qiskit Fall Fest Hackathon 2026 --- Track 2: Quantum Optimization}
}

\date{\today}

\begin{document}

\maketitle

\begin{abstract}
Constrained combinatorial portfolio selection on noisy intermediate-scale quantum (NISQ) hardware is traditionally impaired by soft-penalty encodings that distort the variational landscape and dilute measurement statistics across exponentially large invalid subspaces. In this work, Team Gizmo introduces a subspace-preserving Quantum Approximate Optimization Algorithm (QAOA) framework for Markowitz mean-variance optimization with hard cardinality constraints. By combining polynomial-depth Dicke state initialization ($|D_K^N\rangle$) with 1D open-boundary $XY$-hopping mixers, the quantum state is strictly confined to the feasible subspace, guaranteeing $100.00\%$ constraint satisfaction across all measurement samples. To overcome the dense all-to-all connectivity of financial covariance matrices on sparse physical topologies, we implement a linear alternating SWAP network that executes all $\binom{N}{2}$ two-qubit interactions in $\mathcal{O}(N)$ depth, unrolling physical wire permutations into the classical measurement register with zero quantum gate overhead. Evaluating an 8-asset S\&P 500 portfolio on Qiskit primitives, Gizmo achieves a $100.00\%$ feasibility ratio and recovers the exact combinatorial ground-state energy ($-0.956407$, Sharpe ratio $2.2222$), while conventional soft-penalty QAOA wastes $60.06\%$ of circuit shots on invalid allocations. Scaling benchmarks demonstrate that Gizmo's sampling density advantage expands exponentially, reaching a $67.6\times$ efficiency gain at $N=20$.
\end{abstract}

\section{Introduction}
Modern portfolio theory \cite{markowitz1952} formulates asset allocation as a trade-off between expected return and variance. When subject to cardinality or budget constraints---such as selecting exactly $K$ assets out of an available universe of $N$ equities---the problem becomes an NP-hard binary quadratic program.

On gate-model quantum processors, such combinatorial optimization problems are conventionally mapped to the Quantum Approximate Optimization Algorithm (QAOA) \cite{farhi2014}. However, standard implementations convert equality constraints into unconstrained quadratic penalty terms $\lambda (\sum_{i} x_i - K)^2$. On NISQ devices, this approach exhibits two catastrophic failure modes:
\begin{enumerate}
    \item \textbf{Penalty Dilemma:} Small penalty coefficients $\lambda$ fail to suppress infeasible allocations, while large $\lambda$ values dominate the cost Hamiltonian, creating steep energy canyons that trap classical optimizers and obscure return-risk trade-offs.
    \item \textbf{Exponential State-Space Dilution:} The dimension of the feasible subspace is $\binom{N}{K}$, whereas the total Hilbert space grows as $2^N$. For $N=16, K=4$, feasible states constitute only $2.78\%$ of the Hilbert space; at $N=20, K=5$, this collapses to $1.48\%$. Standard transverse-field mixers scatter quantum amplitudes indiscriminately across invalid states, wasting $>98\%$ of the experimental shot budget.
\end{enumerate}

Project Gizmo resolves these bottlenecks through a hardware-aware, subspace-confined quantum pipeline designed specifically for IBM Quantum heavy-hex architectures.

\section{Problem Formulation \& Ising Mapping}

\subsection{The Constrained Markowitz Model}
Let $N$ denote the number of candidate assets, $\boldsymbol{\mu} \in \mathbb{R}^N$ the vector of annualized expected returns, and $\boldsymbol{\Sigma} \in \mathbb{R}^{N \times N}$ the annualized asset return covariance matrix. For a binary selection vector $\mathbf{x} \in \{0, 1\}^N$ and risk-aversion parameter $q > 0$, the cardinality-constrained Markowitz problem is:
\begin{equation}
\min_{\mathbf{x} \in \{0, 1\}^N} \mathcal{C}(\mathbf{x}) = \frac{q}{2} \mathbf{x}^T \boldsymbol{\Sigma} \mathbf{x} - \boldsymbol{\mu}^T \mathbf{x}
\end{equation}
subject to the exact budget constraint:
\begin{equation}
\sum_{i=0}^{N-1} x_i = K
\end{equation}

\subsection{Mapping to Pauli Spin Operators}
We map binary variables $x_i \in \{0, 1\}$ to spin observables via $x_i = \frac{I - Z_i}{2}$. Expanding the quadratic objective yields the cost Hamiltonian:
\begin{equation}
H_C = \sum_{i=0}^{N-1} h_i Z_i + \sum_{0 \le i < j < N} J_{ij} Z_i Z_j + C_{\text{offset}} I
\end{equation}
where the coupling coefficients are given analytically by:
\begin{align}
h_i &= \frac{1}{2}\mu_i - \frac{q}{4}\Sigma_{ii} - \frac{q}{4} \sum_{j \neq i} \Sigma_{ij} \\
J_{ij} &= \frac{q}{4} \Sigma_{ij} \\
C_{\text{offset}} &= \frac{1}{2}\sum_{i=0}^{N-1} \left(\frac{q}{2}\Sigma_{ii} - \mu_i\right) + \sum_{0 \le i < j < N} \frac{q}{4}\Sigma_{ij}
\end{align}

\begin{table}[t]
\centering
\caption{Asset Universe Characteristics ($N=8$, Historical 1-Year S\&P 500 Daily Returns)}
\label{tab:assets}
\vspace{2pt}
\resizebox{\columnwidth}{!}{%
\begin{tabular}{llcc}
\toprule
\textbf{Ticker} & \textbf{Sector} & \textbf{Expected Return} & \textbf{Volatility} \\
\midrule
AAPL & Technology & 31.42\% & 18.25\% \\
MSFT & Technology & 24.18\% & 16.90\% \\
NVDA & Semiconductors & 78.50\% & 38.12\% \\
JPM  & Financials & 28.60\% & 15.40\% \\
XOM  & Energy & 18.90\% & 16.10\% \\
JNJ  & Healthcare & 12.40\% & 11.20\% \\
PG   & Consumer Staples & 14.10\% & 12.80\% \\
AMZN & Consumer Discr. & 36.20\% & 22.40\% \\
\bottomrule
\end{tabular}%
}
\end{table}

\section{Subspace-Preserving Variational Optimization}

\subsection{Dicke State Initialization}
Rather than preparing the unconstrained uniform superposition $|+\rangle^{\otimes N} = 2^{-N/2} \sum_{\mathbf{x}} |\mathbf{x}\rangle$, Gizmo initializes the quantum register directly in the Dicke state:
\begin{equation}
|D_K^N\rangle = \binom{N}{K}^{-1/2} \sum_{\substack{\mathbf{x} \in \{0, 1\}^N \\ \text{wt}(\mathbf{x}) = K}} |\mathbf{x}\rangle
\end{equation}
To avoid the exponential classical memory overhead of generic statevector decomposition, we implement a deterministic gate ladder based on controlled Givens rotations \cite{bartschi2019}. The circuit initializes $K$ physical qubits in the state $|1\rangle$ and sequentially distributes amplitudes via two-qubit Givens blocks:
\begin{equation}
\setlength{\arraycolsep}{2.5pt}
G(\theta) = \begin{pmatrix} 1 & 0 & 0 & 0 \\ 0 & \cos\theta & \sin\theta & 0 \\ 0 & -\sin\theta & \cos\theta & 0 \\ 0 & 0 & 0 & 1 \end{pmatrix}
\end{equation}
where $\theta = \arccos(\sqrt{j / l})$ for inductive indices $l, j$. Each block is synthesized natively with two CNOT gates and one single-qubit controlled-$R_y$ rotation, requiring $\mathcal{O}(NK)$ total two-qubit gates.

\subsection{The \texorpdfstring{$XY$}{XY} Hopping Mixer}
To preserve total excitation number throughout variational evolution, Gizmo replaces the standard transverse-field mixer $H_M = \sum_i X_i$ with the $XY$ hopping Hamiltonian:
\begin{equation}
H_{XY} = \frac{1}{2} \sum_{(i, j) \in E} \left( X_i X_j + Y_i Y_j \right)
\end{equation}

\begin{theorem}[Subspace Invariance]
Let $Q = \sum_{i=0}^{N-1} \frac{I - Z_i}{2}$ be the total Hamming weight operator. Then $[H_{XY}, Q] = 0$, and the unitary evolution operator $U_{XY}(\beta) = e^{-i \beta H_{XY}}$ maps any state within the $K$-excitation subspace strictly back into the $K$-excitation subspace.
\end{theorem}
\begin{proof}
Consider a single interacting edge $(i, j)$. Expanding the commutator:
\begin{align*}
[X_i X_j + Y_i Y_j, Z_i + Z_j] &= [X_i X_j, Z_i] + [X_i X_j, Z_j] \\
&\quad + [Y_i Y_j, Z_i] + [Y_i Y_j, Z_j]
\end{align*}
Using $[X, Z] = -2iY$ and $[Y, Z] = 2iX$:
\begin{align*}
[X_i X_j, Z_i] &= -2i Y_i X_j, \quad [X_i X_j, Z_j] = -2i X_i Y_j \\
[Y_i Y_j, Z_i] &= 2i X_i Y_j, \quad [Y_i Y_j, Z_j] = 2i Y_i X_j
\end{align*}
Summing these terms yields:
\begin{equation}
[X_i X_j + Y_i Y_j, Z_i + Z_j] = 0
\end{equation}
Because terms on disjoint qubits commute trivially, $[H_{XY}, Q] = 0$. By Stone's theorem, $[e^{-i \beta H_{XY}}, Q] = 0$.
\end{proof}

\subsection{1D Nearest-Neighbor Topology}
To ensure strict physical compatibility with IBM heavy-hex processors, we define the interaction graph $E$ as an open 1D chain:
\begin{equation}
E_{\text{chain}} = \{(i, i+1) \mid 0 \le i \le N-2\}
\end{equation}
This eliminates cross-chip routing required by ring topologies and permits parallel execution across even pairs and odd pairs in exactly two gate layers.

\begin{table*}[t]
\centering
\caption{Performance Comparison: Combinatorial Benchmark ($N=8, K=3$, S\&P 500 Equities, 4,096 Total Shots)}
\label{tab:results}
\vspace{2pt}
\resizebox{\textwidth}{!}{%
\begin{tabular}{lccccc}
\toprule
\textbf{Optimization Approach} & \textbf{Feasible Shots} & \textbf{Feasibility Ratio} & \textbf{Min Energy Found} & \textbf{Exact Optima Hit} & \textbf{Circuit Depth} \\
\midrule
Classical Brute Force (Ground Truth) & 56 (Total) & 100.00\% & -0.956407 & Yes (Global) & --- \\
Naive QAOA (Penalty $\lambda = 2.5$) & 1,636 & 39.94\% & -0.956407 & Yes & 15 \\
\textbf{Gizmo Subspace QAOA} & \textbf{4,096} & \textbf{100.00\%} & \textbf{-0.956407} & \textbf{Yes} & 20 \\
\textbf{Gizmo SWAP-Net (1D Chain)} & \textbf{4,096} & \textbf{100.00\%} & \textbf{-0.956407} & \textbf{Yes} & \textbf{22 (Bounded $\mathcal{O}(N)$)} \\
\textbf{Gizmo COBYLA Variational Loop} & \textbf{4,096} & \textbf{100.00\%} & \textbf{-0.956407} & \textbf{Yes} & 22 \\
\bottomrule
\end{tabular}%
}
\end{table*}

\section{Hardware-Aware SWAP Networks}

\subsection{The All-to-All Connectivity Bottleneck}
The empirical covariance matrix $\boldsymbol{\Sigma}$ is dense, resulting in $\binom{N}{2}$ pairwise $Z_i Z_j$ interaction terms in $H_C$. Mapping a complete interaction graph $K_N$ onto a linear physical coupling topology via standard heuristic routing introduces dozens of uncontrolled SWAP operations that rapidly exceed qubit coherence times.

\subsection{Linear Alternating SWAP Execution}
Gizmo implements a linear alternating SWAP network \cite{kivlichan2018}. Over $N$ sequential rounds, alternating nearest-neighbor pairs apply:
\begin{enumerate}
    \item A parameterized two-qubit rotation $R_{ZZ}(2 \gamma J_{q_a, q_b})$ between physical wires $p_1, p_2$ holding logical assets $q_a, q_b$.
    \item A physical SWAP gate exchanging the logical register assignments.
\end{enumerate}
Over exactly $N$ stages, all $\frac{N(N-1)}{2}$ interactions are executed strictly between physically adjacent qubits in total depth $\mathcal{O}(N)$.

\subsection{Permutation Unrolling at Measurement}
At the completion of the SWAP network, the logical register is permuted: physical wire $p$ holds logical asset $\sigma(p)$. Rather than appending an inverse SWAP network (which would double the circuit depth), Gizmo maps the permutation directly into the classical measurement register:
\begin{equation}
\text{Measure: Wire } p \longrightarrow \text{Classical Bit } \sigma(p)
\end{equation}
This unrolls the permutation with \textbf{zero quantum gate overhead}, ensuring measured bitstrings strictly match logical asset indices.

\section{Empirical Evaluation \& Results}

\subsection{Experimental Setup}
We evaluate the framework using Qiskit 2.x primitives (`StatevectorSampler`) on an 8-asset universe with budget constraint $K=3$ and risk-aversion $q=1.0$. Classical ground truth was established by evaluating all $\binom{8}{3} = 56$ candidate portfolios.

\subsection{Classical Combinatorial Ground Truth}
Exhaustive evaluation identified the optimal portfolio:
\begin{equation}
\mathbf{x}^* = \{\text{AAPL}, \text{XOM}, \text{JNJ}\} \quad (\text{Bitstring: } 00110001)
\end{equation}
with expected annual return of $\mathbf{34.99\%}$, annualized volatility of $\mathbf{14.40\%}$, an optimal Sharpe ratio of $\mathbf{2.2222}$ (relative to $r_f = 3.0\%$), and ground-state cost $\mathbf{E_{\min} = -0.956407}$. The worst-case feasible portfolio in the subspace yielded $E_{\max} = 0.179200$.

\subsection{Comparative QAOA Benchmarks}
As documented in Table~\ref{tab:results}, standard QAOA with penalty parameter $\lambda = 2.5$ produced $1,636$ feasible shots out of $4,096$, yielding a feasibility ratio of only $\mathbf{39.94\%}$. Over $\mathbf{60.06\%}$ of measurement resources were discarded on illegal allocations ($\sum x_i \neq 3$).

In contrast, Gizmo Subspace QAOA, Gizmo SWAP-Net, and the COBYLA-optimized variational loop recorded $\mathbf{4,096 / 4,096}$ valid samples, achieving a verified $\mathbf{100.00\%}$ feasibility ratio. Furthermore, the linear SWAP network compiled to depth 22 while correctly recovering the exact combinatorial global minimum ($-0.956407$).

\begin{table}[b]
\centering
\caption{Combinatorial Scaling: Feasible Subspace Ratio vs. Problem Size ($K = \lfloor N/4 \rfloor$)}
\label{tab:scaling}
\vspace{2pt}
\resizebox{\columnwidth}{!}{%
\begin{tabular}{cccc}
\toprule
\textbf{Assets ($N$)} & \textbf{Budget ($K$)} & \textbf{Naive Space $\frac{\binom{N}{K}}{2^N}$} & \textbf{Gizmo Subspace} \\
\midrule
6  & 2 & 23.44\% & \textbf{100.00\%} \\
8  & 2 & 10.94\% & \textbf{100.00\%} \\
10 & 2 & 4.39\%  & \textbf{100.00\%} \\
12 & 3 & 5.37\%  & \textbf{100.00\%} \\
14 & 3 & 2.22\%  & \textbf{100.00\%} \\
16 & 4 & 2.78\%  & \textbf{100.00\%} \\
20 & 5 & 1.48\%  & \textbf{100.00\%} \\
\bottomrule
\end{tabular}%
}
\end{table}

\subsection{The Exponential Scaling Wall}
Table~\ref{tab:scaling} reports the theoretical and empirical decay of the unconstrained feasible space as $N$ scales from 6 to 20. At $N=20, K=5$, conventional QAOA retains only $1.4786\%$ valid states. By strictly confining evolution to $\binom{N}{K}$, Gizmo provides an empirical sampling concentration advantage of:
\begin{equation}
\text{Advantage Factor} = \frac{2^N}{\binom{N}{K}} \approx 67.6\times \quad (\text{at } N=20)
\end{equation}

\section{Epistemic Limits \& The HUBO Horizon}
Intellectual honesty requires delineating where classical algorithms dominate versus where quantum advantage can realistically reside.

For standard binary quadratic programs with $N \le 50$, classical mixed-integer solvers (e.g., Gurobi, CPLEX, SCIP) reach global optimality in milliseconds via branch-and-cut and polyhedral relaxations. We make no assertion of quantum advantage over classical MIPs on toy quadratic Markowitz instances.

Rather, Gizmo's subspace-preserving architecture provides the essential foundation for \textbf{Higher-Order Binary Optimization (HUBO)}, where classical linear/quadratic relaxations break down:
\begin{itemize}
    \item \textbf{Nonlinear Transaction Friction:} Piecewise cubic market-impact functions ($c_i x_i^{3/2}$).
    \item \textbf{High-Order Risk Tensors:} Asymmetry and tail-risk optimization incorporating co-skewness tensors ($S_{ijk} x_i x_j x_k$) and co-kurtosis tensors ($K_{ijkl} x_i x_j x_k x_l$).
\end{itemize}
Under HUBO formulations, classical branch-and-bound exhibits worst-case exponential tree exploration, whereas multi-qubit phase rotations natively encode multi-body Pauli interactions.

\section{Conclusion}
Project Gizmo demonstrates that hard cardinality constraints in quantum financial optimization can be solved without heuristic penalty parameters. By unifying deterministic Dicke state preparation, 1D nearest-neighbor $XY$-hopping mixers, and linear alternating SWAP networks with measurement unrolling, Gizmo guarantees $100\%$ constraint satisfaction while preserving $\mathcal{O}(N)$ circuit depth on physical IBM heavy-hex hardware.

\begin{thebibliography}{9}
\bibitem{markowitz1952}
H.~Markowitz, ``Portfolio selection,'' \emph{The Journal of Finance}, vol.~7, no.~1, pp.~77--91, 1952.

\bibitem{farhi2014}
E.~Farhi, J.~Goldstone, and S.~Gutmann, ``A quantum approximate optimization algorithm,'' \emph{arXiv:1411.4028}, 2014.

\bibitem{bartschi2019}
A.~B{\"a}rtschi and S.~Eidenbenz, ``Deterministic preparation of Dicke states,'' \emph{Fundamentals of Computation Theory}, pp.~126--139, Springer, 2019.

\bibitem{kivlichan2018}
I.~D.~Kivlichan et~al., ``Quantum simulation of electronic structure with linear depth and connectivity,'' \emph{Physical Review Letters}, vol.~120, no.~11, p.~110501, 2018.
\end{thebibliography}

\end{document}
